\chapter{Introduction}
\label{chap:intro}
\chaptermark{Introduction}

Modern cloud-based software systems are increasingly expected to be reliable, scalable, and highly available. As organizations adopt cloud-native practices and deploy distributed applications, the need for platforms that enable continuous monitoring, maintenance, and delivery of production systems becomes essential. The majority of organizations rely on tools such as Kubernetes \cite{k8s_docs} or HashiCorp Nomad \cite{nomad_docs}; however, these platforms often impose high operational costs and require substantial technical expertise to maintain. For example, deploying a Kubernetes cluster requires at least 4 CPU cores and 16 gigabytes of RAM for the control-plane component alone, and even then, managing upgrades, configuring networking and storage, handling security policies, and troubleshooting failures remain highly complex tasks that often demand dedicated DevOps teams and extensive ongoing maintenance.

The \textit{orchestration principle} \cite{sebrechts2018orchestrator, tosca2015simple} underlies the aforementioned platforms. Recent advances in self-healing systems \cite{rauba2024neurips} and zero-trust security models \cite{vankayalapati2025zerotrust} provide additional mechanisms for enhancing system reliability and security. Orchestration refers to the automated management, coordination, and deployment of multiple services and components in distributed systems to ensure their consistent and efficient operation. However, modern orchestration platforms are often difficult to modify and extend. For instance, platforms that rely on a push-based scheduling model are susceptible to issues such as delivery failure following network connectivity errors and workload duplication. For businesses, these issues can be critical, since infrastructure failures may result in significant financial losses.

The complexity of existing orchestration platforms, combined with the challenges of maintainability and workload management, creates a significant barrier for organizations that cannot afford large DevOps departments. Reliance on a push-based scheduling model limits fault tolerance and introduces bottlenecks in distributed coordination. This thesis addresses these limitations by exploring a pull-based approach to distributed workload management. In the pull-based model, rather than directing a worker node to perform a specific task, the task is placed in a shared registry from which worker nodes independently retrieve and execute it. In practice, instead of receiving instructions from a central controller, worker nodes periodically poll the control plane for available workloads and autonomously claim ownership through a lease mechanism. This model aims to increase reliability, simplify maintenance, and make distributed orchestration more accessible.

The purpose of this work is to design and implement a working prototype of a distributed workload management system based on a \textit{pull-model approach}. The key objectives are as follows: reducing operational overhead; designing a decentralized architecture with the ability to scale; improving security by automating manual infrastructure operations; and creating a lightweight, flexible alternative to existing orchestration platforms. During each phase, artifacts will be collected to ensure the solution can be replicated.

To achieve the stated objectives, the Design Science Research \cite{dresch2014design} methodology is applied, which focuses on building and evaluating artifacts to solve real-world problems. This approach emphasizes iterative development, where each stage informs and refines the next, ensuring that the resulting system effectively addresses the identified challenges. The work is organized into three stages:

\begin{enumerate}
\item Problem identification and requirements clarification: analyzing existing orchestration solutions, identifying their limitations, and defining the functional and non-functional requirements for the proposed system.
\item System design and prototype development: developing the architecture and implementing a working prototype of the distributed workload management system, incorporating features such as task scheduling, fault tolerance, and scalability.
\item Validation and evaluation: testing the prototype under controlled and realistic conditions, measuring its performance, reliability, and resource efficiency, and comparing it against existing platforms to assess improvements and identify areas for further refinement.
\end{enumerate}

This methodology ensures a structured, evidence-based approach to designing and assessing the proposed system while keeping practical applicability as a central focus.

Generative artificial intelligence (GenAI) tools were used in this research as a language and editing aid. GenAI assisted in refining phrasing, improving grammatical correctness, and enhancing the academic tone of the exposition.

The expected outcome of this research is a lightweight and reliable orchestration system prototype that demonstrates the practical benefits of applying a pull-based model in distributed workload management. The developed system is expected to reduce operational overhead by automating configuration delivery and minimizing manual interventions. It will introduce a decentralized control plane that relies on REST API and PostgreSQL for efficient coordination and data storage. Additionally, the work aims to provide a clear comparison with existing orchestration platforms, highlighting improvements in scalability, fault tolerance, and resource efficiency. As a result, this thesis contributes to simplifying infrastructure management and offers an alternative orchestration approach that can be more easily adopted by organizations with different operational scales.

\section{Research Questions}

This thesis seeks to answer the following research questions:

\begin{enumerate}[label=\textbf{RQ\arabic*.}, left=0pt]
\item How can a pull-based orchestration approach reduce operational complexity and resource overhead compared to traditional push-based platforms like Kubernetes, while maintaining comparable fault tolerance and self-healing capabilities?

\item What are the architectural trade-offs inherent in a fully pull-based orchestration model, particularly with respect to consistency guarantees, coordination overhead, fault isolation mechanisms, and the trade-off between scheduling simplicity and placement optimality?
\end{enumerate}

\section{Thesis Structure}

The remainder of this thesis is organized as follows:

\begin{description}[style=multiline, leftmargin=3cm, itemsep=0.8em]
\item[\textbf{Chapter 1: Introduction}] Motivates the problem, states the research questions, outlines the Design Science Research methodology, and provides an overview of the thesis structure.

\item[\textbf{Chapter 2: Literature Review}] Investigate modern container orchestration systems and decentralized, pull-based approaches in distributed computing, highlighting gaps in current research and practice.

\item[\textbf{Chapter 3: Problem Identification}] Analyzes the business impact of orchestration failures, examines existing solutions (Kubernetes, HashiCorp Nomad), and derives functional and non-functional requirements for the proposed system.

\item[\textbf{Chapter 4: Implementation}] Presents the system architecture, design decisions, and technical implementation details through key scenarios, demonstrating how the pull-based model operates in practice.

\item[\textbf{Chapter 5: Evaluation and Discussion}] Evaluates the prototype against the established requirements, measures performance metrics (including RTO achievement), and discusses the architectural trade-offs observed in the pull-based model.

\item[\textbf{Chapter 6: Conclusion}] Summarizes the contributions, limitations, and practical implications of the research, and outlines directions for future work.
\end{description}
